\lyseis{Γ3.3}{Κλασικός ορισμός της πιθανότητας}
\sol{1} Ο δειγματικός χώρος έχει $6$ ισοπίθανα στοιχεία. α) $\dfrac{1}{6}\approx16{,}7\%$. β) $\{2,4,6\}$: $\dfrac{3}{6}=\dfrac{1}{2}$. γ) $\{1,2\}$: $\dfrac{2}{6}=\dfrac{1}{3}$. δ) $\{1,2,3,4,5\}$: $\dfrac{5}{6}$.
\sol{2} Οι δυνατές μπάλες είναι $8$. α) $\dfrac{4}{8}=\dfrac{1}{2}=0{,}5=50\%$. β) $\dfrac{3}{8}=0{,}375=37{,}5\%$. γ) Όχι κίτρινη: $7$ μπάλες, $\dfrac{7}{8}=0{,}875=87{,}5\%$. δ) Δεν υπάρχει πράσινη: $\dfrac{0}{8}=0=0\%$.
\sol{3} $\frac{1}{2}=0{,}5=50\%$. $0{,}25=\frac{1}{4}=25\%$. $30\%=\frac{3}{10}=0{,}3$. $\frac{3}{5}=0{,}6=60\%$. $0{,}05=\frac{1}{20}=5\%$. $100\%=1=\frac{1}{1}$.
\sol{4} α) Λάθος: δεν ξεπερνά το $1$. β) Σωστό. γ) Σωστό: $1-0{,}3=0{,}7$. δ) Σωστό. ε) Λάθος: δεν υπάρχει το $7$, άρα η πιθανότητα είναι $0$. στ) Λάθος: η πιθανότητα δεν εγγυάται ακριβή αριθμό φορών.
\sol{5} α) Βέβαιο. β) Αδύνατο. γ) Ούτε βέβαιο ούτε αδύνατο ($\frac{1}{2}$). δ) Αδύνατο. ε) Βέβαιο.
\sol{6} Η λέξη έχει $10$ γράμματα (Μ, Α, Θ, Η, Μ, Α, Τ, Ι, Κ, Α). α) Το Α υπάρχει $3$ φορές: $\dfrac{3}{10}=0{,}3$. β) Το Μ υπάρχει $2$ φορές: $\dfrac{2}{10}=\dfrac{1}{5}=0{,}2$. γ) Φωνήεντα (Α, Η, Α, Ι, Α): $5$ γράμματα, $\dfrac{5}{10}=\dfrac{1}{2}$. δ) Δεν υπάρχει το Ω: $0$.
\sol{7} Ισοπίθανα είναι τα $36$ ζευγάρια. α) Άθροισμα $8$: $(2,6),(3,5),(4,4),(5,3),(6,2)$, $5$ ζευγάρια: $\dfrac{5}{36}\approx13{,}9\%$. β) $(1,1),\dots,(6,6)$: $\dfrac{6}{36}=\dfrac{1}{6}$. γ) Άθροισμα $1$ δεν υπάρχει: $\dfrac{0}{36}=0$.
\sol{8} Κερδίζει: $\dfrac{10}{200}=\dfrac{1}{20}=0{,}05=5\%$. Δεν κερδίζει: $\dfrac{190}{200}=\dfrac{19}{20}=0{,}95=95\%$.
\sol{9} Κορίτσι: $\dfrac{14}{25}=0{,}56=56\%$. Αγόρι: $\dfrac{11}{25}=0{,}44=44\%$.
\sol{10} Κόκκινο: $\dfrac{5}{12}\approx41{,}7\%$. Μπλε: $\dfrac{4}{12}=\dfrac{1}{3}\approx33{,}3\%$. Πράσινο: $\dfrac{2}{12}=\dfrac{1}{6}\approx16{,}7\%$. Χρυσό: $\dfrac{1}{12}\approx8{,}3\%$. Άθροισμα $\dfrac{5+4+2+1}{12}=\dfrac{12}{12}=1$, γιατί σε κάθε στροφή ο τροχός σταματά σίγουρα σε κάποιο χρώμα (βέβαιο ενδεχόμενο).
\sol{11} Υπάρχουν $10$ ψηφία και τα περιττά είναι $1,3,5,7,9$, δηλαδή $5$. Αφού ξέρει ότι είναι περιττό, μαντεύει ανάμεσα σε $5$ ισοπίθανα: $\dfrac{1}{5}=0{,}2=20\%$.
\sol{12} α) Οι κόκκινοι άσοι είναι $2$ (κούπα, καρό): $\dfrac{2}{52}=\dfrac{1}{26}\approx3{,}8\%$. β) Φιγούρες $12$: $\dfrac{12}{52}=\dfrac{3}{13}\approx23{,}1\%$. γ) Όχι φιγούρα: $52-12=40$, $\dfrac{40}{52}=\dfrac{10}{13}\approx76{,}9\%$ (και $1-\frac{3}{13}=\frac{10}{13}$).
\sol{13} Κόκκινες: $0{,}35\cdot20=7$. Μπλε: $\dfrac{2}{5}\cdot20=8$. Πράσινες: $20-7-8=5$.
\sol{14} Ο δειγματικός χώρος έχει $8$ ισοπίθανα στοιχεία $\mathrm{ΚΚΚ},\ \mathrm{ΚΚΓ},\ \mathrm{ΚΓΚ},\ \mathrm{ΚΓΓ},\ \mathrm{ΓΚΚ},\ \mathrm{ΓΚΓ},\ \mathrm{ΓΓΚ},\ \mathrm{ΓΓΓ}$. α) Μόνο το $\mathrm{ΚΚΚ}$: $\dfrac{1}{8}=0{,}125=12{,}5\%$. β) Όλα εκτός από το $\mathrm{ΚΚΚ}$: $\dfrac{7}{8}=0{,}875=87{,}5\%$. γ) $\mathrm{ΚΚΓ},\ \mathrm{ΚΓΚ},\ \mathrm{ΓΚΚ}$: $\dfrac{3}{8}=0{,}375=37{,}5\%$.
\sol{15} Όχι. Τα δύο αποτελέσματα «βρέχει» και «δεν βρέχει» δεν είναι ισοπίθανα, οπότε δεν μπορούμε να χρησιμοποιήσουμε τον κλασικό ορισμό. Στο νόμισμα τα δύο αποτελέσματα είναι ισοπίθανα λόγω της συμμετρίας του νομίσματος. Η πιθανότητα βροχής εκτιμάται από πολλές παρόμοιες μέρες στο παρελθόν (σχετική συχνότητα).
\sol{16} $\dfrac{5}{20}=0{,}25$, $\dfrac{18}{100}=0{,}18$, $\dfrac{98}{600}\approx0{,}163$. Η πιθανότητα είναι $\frac{1}{6}\approx0{,}167$. Με $20$ ρίψεις η σχετική συχνότητα απέχει αρκετά, με $600$ είναι πολύ κοντά. Όσο αυξάνονται οι ρίψεις, η σχετική συχνότητα πλησιάζει την πιθανότητα.
\sol{17} Ελεύθερη κατασκευή. Χρειάζονται $3$ κόκκινα τμήματα, $\frac{1}{4}=\frac{2}{8}$, άρα $2$ μπλε, και τα υπόλοιπα $8-3-2=3$ τμήματα έχουν άλλα χρώματα.
\sol{18} Τα $11$ αθροίσματα δεν είναι ισοπίθανα. Ισοπίθανα είναι τα $36$ ζευγάρια. Το $7$ προκύπτει από $6$ ζευγάρια, άρα η πιθανότητα είναι $\dfrac{6}{36}=\dfrac{1}{6}$, όχι $\dfrac{1}{11}$.
\sol{19} Ελεύθερη εργασία. Αν π.χ. ήρθε Κεφάλι $9$ φορές, η σχετική συχνότητα είναι $\frac{9}{20}=0{,}45$, κοντά στο $0{,}5$ αλλά όχι ακριβώς. Με άλλους φίλους και περισσότερες ρίψεις θα πλησιάζει περισσότερο.
\sol{20} Ελεύθερη εργασία. Παράδειγμα: κουτί με $3$ κόκκινες, $5$ μπλε και $2$ πράσινες μπάλες. Η πιθανότητα για μπλε είναι $\dfrac{5}{10}=\dfrac{1}{2}=0{,}5=50\%$.
